\documentclass[11pt,a4paper]{article}

\usepackage[margin=2.5cm]{geometry}
\usepackage{amsmath,amssymb}
\usepackage{booktabs}
\usepackage{hyperref}

\newcommand{\vect}[1]{\mathbf{#1}}
\newcommand{\erfc}{\operatorname{erfc}}
\newcommand{\um}{\,\mu\mathrm{m}}

\title{Ewald summation of the chemoattractant field\\ of a periodic phytoplankton community}
\author{}
\date{}

\begin{document}
\maketitle

\section{Chemoattractant field}

Each phytoplankton cell $j$ is a sphere of radius $R_j$ centred at $\vect{c}_j$ that releases
chemoattractant at a constant rate. The chemoattractant diffuses and is consumed at a constant
rate, which defines a decay length $\gamma$. At steady state, the concentration outside an isolated
cell is
\begin{equation}
  C_j(r) = C_{s,j}\,\frac{R_j}{r}\,e^{-(r-R_j)/\gamma}, \qquad r = |\vect{x}-\vect{c}_j| \ge R_j,
  \label{eq:single}
\end{equation}
where $C_{s,j}$ is the concentration at the cell surface. Outside the cell, Eq.~\eqref{eq:single}
is the field of a point source at $\vect{c}_j$,
\begin{equation}
  C_j(r) = Q_j\,G(r), \qquad G(r) = \frac{e^{-r/\gamma}}{r}, \qquad Q_j = C_{s,j}\,R_j\,e^{R_j/\gamma},
  \label{eq:pointsource}
\end{equation}
where $G$ is the screened Coulomb (Yukawa) kernel. This representation is exact for $r \ge R_j$,
which is the only region where the field is needed, since bacteria cannot enter the cells.

The community is a set of $N_s$ cells in a cubic box of side $L$ with periodic boundary conditions.
The concentration at a point $\vect{x}$ in the fluid is the sum of the contributions of all cells
and of all their periodic images,
\begin{equation}
  C(\vect{x}) = C_b + \sum_{j=1}^{N_s} \sum_{\vect{n}\in\mathbb{Z}^3} Q_j\,
  G\bigl(|\vect{x} - \vect{c}_j - \vect{n}L|\bigr),
  \label{eq:fullsum}
\end{equation}
where $C_b$ is a uniform background. The sum converges absolutely because $G$ decays exponentially,
but the number of terms that contribute is large when $\gamma$ is comparable to or larger than the
distance between cells. Each bacterium needs $C$ and $\nabla C$ at its position at every time step,
so we evaluate Eq.~\eqref{eq:fullsum} by Ewald summation.

\section{Ewald splitting of the screened kernel}

We split the kernel into a short-range and a long-range part,
$G = G_{\mathrm{near}} + G_{\mathrm{far}}$, using the standard Ewald splitting for Yukawa
interactions~\cite{Ewald1921,SalinCaillol2000}. With $\kappa = 1/\gamma$ and a splitting length
$\sigma$, the long-range part is defined by its Fourier transform,
\begin{equation}
  \hat G_{\mathrm{far}}(k) = \frac{4\pi}{k^2+\kappa^2}\,e^{-(k^2+\kappa^2)\sigma^2/4},
  \label{eq:gfar_k}
\end{equation}
which is the Fourier transform of $G$, $4\pi/(k^2+\kappa^2)$, multiplied by a Gaussian. In real
space,
\begin{align}
  G_{\mathrm{far}}(r) &= \frac{1}{2r}\Bigl[e^{-\kappa r}\erfc\!\Bigl(\frac{\kappa\sigma}{2} - \frac{r}{\sigma}\Bigr)
    - e^{\kappa r}\erfc\!\Bigl(\frac{\kappa\sigma}{2} + \frac{r}{\sigma}\Bigr)\Bigr], \\
  G_{\mathrm{near}}(r) &= \frac{1}{2r}\Bigl[e^{-\kappa r}\erfc\!\Bigl(\frac{r}{\sigma} - \frac{\kappa\sigma}{2}\Bigr)
    + e^{\kappa r}\erfc\!\Bigl(\frac{r}{\sigma} + \frac{\kappa\sigma}{2}\Bigr)\Bigr].
  \label{eq:gnear}
\end{align}
$G_{\mathrm{far}}$ is smooth and finite at $r=0$ and tends to $G$ for $r \gg \sigma$.
$G_{\mathrm{near}}$ carries the $1/r$ singularity and decays like $\erfc(r/\sigma)$, i.e.\ it is
negligible beyond a few $\sigma$. The Gaussian factor in Eq.~\eqref{eq:gfar_k} contains
$k^2+\kappa^2$ rather than $k^2$ alone; this is what makes $G_{\mathrm{far}} \to G$ at large
distance, so that $G_{\mathrm{near}}$ vanishes there. The radial derivative used for the gradient is
\begin{equation}
  \frac{dG_{\mathrm{near}}}{dr} = \frac{1}{2r}\Bigl[\kappa\,e^{\kappa r}\erfc(u+a) - \kappa\,e^{-\kappa r}\erfc(u-a)
    - \frac{4}{\sigma\sqrt{\pi}}\,e^{-u^2-a^2}\Bigr] - \frac{G_{\mathrm{near}}(r)}{r},
\end{equation}
with $u = r/\sigma$ and $a = \kappa\sigma/2$.

The concentration is then $C = C_b + C_{\mathrm{near}} + C_{\mathrm{far}}$, where each part is
Eq.~\eqref{eq:fullsum} with $G$ replaced by $G_{\mathrm{near}}$ or $G_{\mathrm{far}}$.
The two parts are computed with different methods.

\section{Short-range part}

The short-range part is summed directly over cells closer than a cutoff $r_c$,
\begin{equation}
  C_{\mathrm{near}}(\vect{x}) = \sum_{j:\,r_j<r_c} Q_j\,G_{\mathrm{near}}(r_j), \qquad
  \nabla C_{\mathrm{near}}(\vect{x}) = \sum_{j:\,r_j<r_c} Q_j\,\frac{dG_{\mathrm{near}}}{dr}(r_j)\,\frac{\vect{r}_j}{r_j},
\end{equation}
where $\vect{r}_j$ is the minimum-image separation between $\vect{x}$ and $\vect{c}_j$, and
$r_j = |\vect{r}_j|$. The pairs are found with a cell list~\cite{CellListMap}. Since
$r_c < L/2$, at most one image of each cell lies within the cutoff, and the minimum image is
sufficient. The cost per step is proportional to the number of bacteria times the mean number of
cells within $r_c$, $\rho_s\,\tfrac{4}{3}\pi r_c^3$, where $\rho_s = N_s/L^3$.

\section{Long-range part}

\paragraph{Fourier representation.}
By the Poisson summation formula, the periodic sum of the long-range part is a Fourier series over
the wave vectors $\vect{k} = 2\pi\vect{n}/L$ of the box,
\begin{equation}
  C_{\mathrm{far}}(\vect{x}) = \frac{1}{L^3}\sum_{\vect{k}} \hat G_{\mathrm{far}}(k)\,S(\vect{k})\,e^{i\vect{k}\cdot\vect{x}},
  \qquad S(\vect{k}) = \sum_{j=1}^{N_s} Q_j\,e^{-i\vect{k}\cdot\vect{c}_j}.
  \label{eq:fourier}
\end{equation}
Because of the Gaussian factor in Eq.~\eqref{eq:gfar_k} the series converges rapidly, and we keep
the $M^3$ wave vectors with $|n_\alpha| < M/2$ in each direction ($M$ even; the Nyquist modes
$n_\alpha = M/2$ are set to zero so that the truncated series is real). The largest neglected
wavenumber is $k_{\max} = \pi M/L = \pi/h$, with $h = L/M$, and the neglected terms are smaller than
the retained ones by a factor of order
\begin{equation}
  \varepsilon_k = e^{-\pi^2\sigma^2/(4h^2)}.
  \label{eq:eps_k}
\end{equation}
The $k=0$ term is finite, $4\pi\gamma^2 e^{-\sigma^2/4\gamma^2} \sum_j Q_j/L^3$, so no
neutrality condition is needed, unlike the Coulomb case.

\paragraph{Structure factor.}
The structure factor $S(\vect{k})$ on the $M^3$ wave vectors is a type-1 non-uniform fast Fourier
transform (NUFFT) of the source strengths $Q_j$ at the positions $\vect{c}_j$~\cite{Barnett2019}.
It is computed with NonuniformFFTs.jl~\cite{NonuniformFFTs}: the sources are spread onto an oversampled grid
of $(2M)^3$ points with a backwards Kaiser--Bessel kernel of half-support $m$, the grid is
transformed by FFT, and the kernel is deconvolved. The use of a NUFFT for the structure factor is
the approach of spectral Ewald methods~\cite{LindboTornberg2011}.

\paragraph{Mesh.}
The Fourier coefficients $\hat G_{\mathrm{far}}(k) S(\vect{k})/L^3$ are transformed to real space
by an inverse real FFT~\cite{FFTW}, which gives $C_{\mathrm{far}}$ at the $M^3$ nodes
$\vect{x}_{\vect{i}} = (\vect{i}-1)\,h$ of a uniform mesh with spacing $h = L/M$. These nodal
values are the truncated series Eq.~\eqref{eq:fourier}, exactly up to the NUFFT tolerance. The
cell positions do not change during a simulation, so this step is done once, when the model is
set up.

\paragraph{Interpolation.}
At every step $C_{\mathrm{far}}$ and $\nabla C_{\mathrm{far}}$ are evaluated at the bacterial
positions by periodic cubic B-spline interpolation of the mesh values~\cite{Unser1993}, as in
smooth particle-mesh Ewald~\cite{Darden1993,Essmann1995}. We use Interpolations.jl~\cite{Interpolations}
(\texttt{BSpline(Cubic(Periodic(OnCell())))}). The B-spline coefficients are obtained once, at
setup, by solving the periodic prefilter system along each axis, so that the interpolant passes
through the mesh values. The \texttt{OnCell} variant is the one whose period is $M$ nodes, which
matches the FFT mesh, where node $M+1$ is node $1$. Each evaluation combines the $4^3$ coefficients
around the point, and the gradient is the analytic derivative of the spline. The interpolation
error of a cubic spline is $O(h^4)$ for the value and $O(h^3)$ for the gradient, relative to the
length scale over which the interpolated function varies. Here that length scale is $\sigma$,
because the Fourier spectrum of $C_{\mathrm{far}}$ is cut off by $e^{-k^2\sigma^2/4}$. The accuracy
of the long-range part is therefore set by the ratio $h/\sigma$.

The cost per step is a fixed number of operations per bacterium, independent of $N_s$, $\sigma$ and
$L$. The setup costs $O(N_s m^3)$ for spreading and $O(M^3\log M)$ for the FFTs, and the mesh uses
$8M^3$ bytes, with $M \approx 2.4\,L/\sigma$ for the parameters below.

\section{Algorithm}

At setup: (i) compute $Q_j$ from Eq.~\eqref{eq:pointsource}; (ii) compute $S(\vect{k})$ by NUFFT;
(iii) multiply by $\hat G_{\mathrm{far}}(k)/L^3$ and transform to the mesh; (iv) compute the
B-spline coefficients. At every time step, after all bacteria have moved: (v) update the cell list
and sum the short-range part over the neighbouring cells; (vi) add the interpolated long-range part
and its gradient. The resulting $C$ and $\nabla C$ are then used by the chemotaxis model of each
bacterium.

\section{Parameters}

\begin{table}[h]
  \centering\small
  \begin{tabular}{llll}
    \toprule
    Parameter & Symbol & Value & Controls \\
    \midrule
    Splitting length & $\sigma$ & $40\um$ & short- vs.\ long-range cost \\
    Short-range cutoff & $r_c$ & $3.5\,\sigma = 140\um$ & truncation of the short-range part \\
    Mesh spacing & $h$ & $\le \sigma/2.4 = 16.7\um$ & interpolation error \\
    Mesh size & $M$ & $2\lceil L/2h\rceil$ & --- \\
    NUFFT kernel half-support & $m$ & 8 & NUFFT tolerance \\
    NUFFT oversampling & & 2 & (NonuniformFFTs.jl default) \\
    \bottomrule
  \end{tabular}
  \caption{Numerical parameters of the Ewald summation. The physical parameters are
    $\gamma = 150\um$ and the community composition.}
  \label{tab:params}
\end{table}

The choices in Table~\ref{tab:params} are based on the tests below. We require a relative error of
at most $10^{-3}$ on $C$ and $10^{-2}$ on $\nabla C$ everywhere in the fluid, and choose the
parameters that meet this requirement at the lowest cost.

\subsection{Accuracy tests}

\paragraph{Test system.} A community of $N_s = 153$ cells (abundance $3\times10^5$\,mL$^{-1}$,
allometric exponent $\alpha=1$, radii $0.65$--$10\um$) in a periodic box of side $L = 800\um$, with
$\gamma = 150\um$. We evaluate the field at 400 points: 200 at random positions $0.5$--$5\um$ from
the surface of a randomly chosen cell, where the field and its gradient are largest, and 200 at
uniformly random positions in the fluid.

\paragraph{Reference.} The reference is the direct sum, Eq.~\eqref{eq:fullsum}, over all cells and
all periodic images closer than $18\gamma = 2700\um$, using the exact single-cell field
Eq.~\eqref{eq:single} and its gradient. Extending the sum to $24\gamma$ changes $C$ by less than
$2\times10^{-7}$ and $\nabla C$ by less than $3\times10^{-10}$ (relative), so the reference is
accurate well beyond the errors below.

\paragraph{Error measures.} For each point we compute the relative error of the concentration,
$|\delta C|/C$, and of the gradient, $|\delta\nabla C|/|\nabla C|$. Far from all cells the
contributions of different cells partly cancel and $|\nabla C|$ can be much smaller than its
typical scale $C/\gamma$. There the relative gradient error is large even when the absolute error is
small, so we also report $|\delta\nabla C|/(C/\gamma)$. Tables~\ref{tab:h}--\ref{tab:sigma} give
the median and the maximum over the 400 points. In each test one parameter is varied and the
others are set to values at which their own contribution to the error is negligible.

\begin{table}[h]
  \centering\small
  \begin{tabular}{lllllll}
    \toprule
    $\sigma/h$ & $h$ ($\mu$m) & $|\delta C|/C$ med. & max & $|\delta\nabla C|/|\nabla C|$ med. & max & $|\delta\nabla C|/(C/\gamma)$ max \\
    \midrule
    1.2 & 33.3 & $4.8\times10^{-4}$ & $4.4\times10^{-3}$ & $1.1\times10^{-3}$ & $6.4\times10^{-2}$ & $3.4\times10^{-2}$ \\
    1.6 & 25.0 & $9.3\times10^{-5}$ & $9.0\times10^{-4}$ & $2.4\times10^{-4}$ & $1.2\times10^{-2}$ & $1.3\times10^{-2}$ \\
    2.0 & 20.0 & $2.9\times10^{-5}$ & $4.1\times10^{-4}$ & $9.8\times10^{-5}$ & $4.8\times10^{-3}$ & $5.1\times10^{-3}$ \\
    \textbf{2.4} & \textbf{16.7} & $1.4\times10^{-5}$ & $9.7\times10^{-5}$ & $5.1\times10^{-5}$ & $3.7\times10^{-3}$ & $2.6\times10^{-3}$ \\
    3.2 & 12.5 & $3.8\times10^{-6}$ & $3.5\times10^{-5}$ & $2.2\times10^{-5}$ & $1.6\times10^{-3}$ & $1.2\times10^{-3}$ \\
    4.0 & 10.0 & $1.5\times10^{-6}$ & $9.6\times10^{-6}$ & $9.8\times10^{-6}$ & $7.8\times10^{-4}$ & $5.4\times10^{-4}$ \\
    4.8 & 8.33 & $7.1\times10^{-7}$ & $5.6\times10^{-6}$ & $5.6\times10^{-6}$ & $3.1\times10^{-4}$ & $3.5\times10^{-4}$ \\
    \bottomrule
  \end{tabular}
  \caption{Error as a function of the mesh spacing ($\sigma = 40\um$, $r_c = 3.5\sigma$, $m = 8$).}
  \label{tab:h}
\end{table}

\paragraph{Mesh spacing.} Table~\ref{tab:h} shows the error as a function of $\sigma/h$. For
$\sigma/h \le 2$ the truncation of the Fourier series contributes: by Eq.~\eqref{eq:eps_k},
$\varepsilon_k = 3\times10^{-2}$, $2\times10^{-3}$ and $5\times10^{-5}$ at $\sigma/h = 1.2$, $1.6$
and $2.0$. From $\sigma/h = 2.4$ ($\varepsilon_k = 7\times10^{-7}$) on, the error is that of the
cubic interpolation. Halving $h$ from $\sigma/2.4$ to $\sigma/4.8$ reduces the maximum error by a
factor 17 on $C$ and 12 on $\nabla C$, consistent with the expected $O(h^4)$ and $O(h^3)$
convergence. We use $h = \sigma/2.4$, the coarsest mesh that meets the accuracy requirement with a
margin: the maximum error is $10^{-4}$ on $C$ and $4\times10^{-3}$ on $\nabla C$. Near the cell
surfaces, where the short-range part dominates, the maximum gradient error is $1.3\times10^{-4}$;
the largest relative gradient errors occur in the bulk, where $|\nabla C| \ll C/\gamma$. A finer
mesh reduces the error further at a memory and setup cost proportional to $h^{-3}$.

\begin{table}[h]
  \centering\small
  \begin{tabular}{lllll}
    \toprule
    $r_c/\sigma$ & $G_{\mathrm{near}}(r_c)/G(r_c)$ & $|\delta C|/C$ max & $|\delta\nabla C|/|\nabla C|$ max & $|\delta\nabla C|/(C/\gamma)$ max \\
    \midrule
    2.0 & $7.9\times10^{-3}$ & $1.5\times10^{-3}$ & $5.1\times10^{-2}$ & $2.5\times10^{-2}$ \\
    2.5 & $7.8\times10^{-4}$ & $2.1\times10^{-4}$ & $1.1\times10^{-2}$ & $3.0\times10^{-3}$ \\
    3.0 & $4.8\times10^{-5}$ & $6.4\times10^{-6}$ & $4.3\times10^{-4}$ & $3.5\times10^{-4}$ \\
    \textbf{3.5} & $1.9\times10^{-6}$ & $5.6\times10^{-6}$ & $3.1\times10^{-4}$ & $3.5\times10^{-4}$ \\
    4.0 & $4.4\times10^{-8}$ & $5.6\times10^{-6}$ & $3.1\times10^{-4}$ & $3.5\times10^{-4}$ \\
    \bottomrule
  \end{tabular}
  \caption{Error as a function of the short-range cutoff ($\sigma = 40\um$, $h = \sigma/4.8$,
    $m = 8$). The second column is the fraction of the kernel that the short-range part still
    carries at the cutoff.}
  \label{tab:rc}
\end{table}

\paragraph{Short-range cutoff.} Table~\ref{tab:rc} uses a fine mesh ($h = \sigma/4.8$), so that the
error from the cutoff is visible. The error stops decreasing at $r_c = 3\sigma$ and the remaining
error is that of the mesh. We use $r_c = 3.5\sigma$, at which the short-range kernel is
$2\times10^{-6}$ of the full kernel, so that the cutoff error is negligible compared to the
interpolation error at $h = \sigma/2.4$. The cutoff also constrains the box size: the minimum-image
convention requires $L > 2r_c = 7\sigma = 280\um$.

\paragraph{NUFFT accuracy.} With $\sigma = 40\um$, $h = \sigma/4.8$ and $r_c = 3.5\sigma$, the
maximum error on $C$ is $2.9\times10^{-4}$, $6.9\times10^{-6}$ and $5.6\times10^{-6}$ for kernel
half-supports $m = 2$, $3$ and $4$, and does not change for $m = 6$ and $8$. The NUFFT error is thus
below the other errors for $m \ge 4$. We use $m = 8$, which makes the NUFFT accurate to near machine
precision. With $N_s \approx 300$ sources the spreading cost, $N_s(2m)^3$ operations, is negligible
compared to the FFTs, so this choice is free.

\begin{table}[h]
  \centering\small
  \begin{tabular}{llllll}
    \toprule
    $\sigma$ ($\mu$m) & $h$ ($\mu$m) & $|\delta C|/C$ max & $|\delta\nabla C|/|\nabla C|$ med. & max & $|\delta\nabla C|/(C/\gamma)$ max \\
    \midrule
    20 & 8.33 & $1.6\times10^{-4}$ & $5.5\times10^{-5}$ & $1.6\times10^{-3}$ & $5.4\times10^{-3}$ \\
    30 & 12.5 & $1.1\times10^{-4}$ & $5.4\times10^{-5}$ & $1.6\times10^{-3}$ & $3.0\times10^{-3}$ \\
    \textbf{40} & 16.7 & $9.7\times10^{-5}$ & $5.1\times10^{-5}$ & $3.7\times10^{-3}$ & $2.6\times10^{-3}$ \\
    60 & 25.0 & $2.0\times10^{-4}$ & $4.5\times10^{-5}$ & $4.7\times10^{-3}$ & $2.3\times10^{-3}$ \\
    80 & 33.3 & $1.4\times10^{-4}$ & $7.0\times10^{-5}$ & $6.8\times10^{-3}$ & $1.9\times10^{-3}$ \\
    100 & 40.0 & $1.6\times10^{-4}$ & $5.3\times10^{-5}$ & $7.6\times10^{-3}$ & $1.2\times10^{-3}$ \\
    \bottomrule
  \end{tabular}
  \caption{Error as a function of the splitting length at fixed $h = \sigma/2.4$, $r_c = 3.5\sigma$,
    $m = 8$.}
  \label{tab:sigma}
\end{table}

\paragraph{Splitting length.} When $h$ and $r_c$ are scaled with $\sigma$, the error depends only
weakly on $\sigma$ (Table~\ref{tab:sigma}): the maximum error on $C$ stays between $1$ and
$2\times10^{-4}$, and the maximum gradient error stays below $10^{-2}$ for all
$20\um \le \sigma \le 100\um$. The choice of $\sigma$ is therefore determined by cost.

\subsection{Cost}

\begin{table}[h]
  \centering\small
  \begin{tabular}{rrrrrrrrr}
    \toprule
    $L$ ($\mu$m) & $\sigma$ ($\mu$m) & $N_s$ & $M$ & mesh (MB) & setup (s) & near (s) & far (s) & full step (s) \\
    \midrule
    1000 & 20 & 298 & 120 & 13.8 & 0.53 & 0.020 & 0.078 & 0.29 \\
    1000 & 30 & 298 & 80 & 4.1 & 0.13 & 0.035 & 0.066 & 0.31 \\
    1000 & \textbf{40} & 298 & 60 & 1.7 & 0.02 & 0.057 & 0.052 & 0.34 \\
    1000 & 60 & 298 & 40 & 0.5 & 0.09 & 0.144 & 0.055 & 0.40 \\
    1000 & 80 & 298 & 30 & 0.2 & $<$0.01 & 0.305 & 0.061 & 0.56 \\
    \midrule
    2150 & 20 & 297 & 258 & 137.4 & 2.0 & 0.011 & 0.127 & 0.39 \\
    2150 & 30 & 297 & 172 & 40.7 & 0.87 & 0.014 & 0.104 & 0.31 \\
    2150 & \textbf{40} & 297 & 130 & 17.6 & 0.27 & 0.019 & 0.092 & 0.30 \\
    2150 & 60 & 297 & 86 & 5.1 & 0.19 & 0.031 & 0.063 & 0.27 \\
    2150 & 80 & 297 & 66 & 2.3 & 0.10 & 0.053 & 0.072 & 0.31 \\
    \bottomrule
  \end{tabular}
  \caption{Cost for $10^5$ bacteria in the two community configurations used in the simulations
    (allometric exponent $\alpha = 0.85$, abundance $3\times10^5$ and $3\times10^4$\,mL$^{-1}$).
    ``Setup'' is the construction of the long-range mesh; ``near'' and ``far'' are the times per
    step of the short- and long-range parts over all bacteria; ``full step'' includes motion,
    collisions and the chemotactic update of all bacteria. Single thread, AMD Ryzen 7 PRO 6850U,
    Julia 1.11.5; minimum of 5 repetitions. Differences of about $0.02$\,s are within the
    run-to-run variation.}
  \label{tab:cost}
\end{table}

Table~\ref{tab:cost} shows the cost for $10^5$ bacteria in the two community configurations used in
the simulations, $\approx 300$ cells in boxes of side $1000\um$ and $2150\um$. The short-range cost
grows as $\sigma^3$, through the number of cells within $r_c$ of a bacterium,
$\rho_s\,\tfrac{4}{3}\pi(3.5\sigma)^3$ (3.4 at $L = 1000\um$, $\sigma = 40\um$; 27 at
$\sigma = 80\um$). The long-range cost per step depends weakly on $\sigma$, through the size of the
mesh and therefore its cache efficiency. Its memory and setup cost grow as $(L/\sigma)^3$.

In the denser configuration ($L = 1000\um$) the field costs about $0.1$\,s per step for
$\sigma \le 40\um$ and becomes dominated by the short-range part for larger $\sigma$. We use
$\sigma = 40\um$, the largest value before the short-range cost increases, because it gives the
smallest mesh. This matters for larger boxes. At fixed $\sigma = 40\um$ the mesh has $M \approx L/16.7\um$ nodes per side, i.e.\
$18$\,MB at $L = 2150\um$ and about $0.2$\,GB at $L = 5000\um$; at $\sigma = 20\um$ these would be
$0.14$ and $1.7$\,GB. At $\sigma = 40\um$ the field takes about one third of the full time step; the
rest is bacterial motion and chemotaxis. The cost per step does not depend on $L$ at fixed $N_s$.

\section{Implementation}

The method is implemented in Julia~1.11.5 in \texttt{src/ewald.jl} (long-range part and
short-range kernel) and \texttt{src/model.jl} (\texttt{community\_field!}), using
NonuniformFFTs.jl~0.9.7~\cite{NonuniformFFTs} for the NUFFT, FFTW.jl~1.10.0~\cite{FFTW} for the
inverse FFT, Interpolations.jl~0.15.2~\cite{Interpolations} for the B-spline interpolation,
SpecialFunctions.jl~2.9.0 for $\erfc$ and CellListMap.jl~0.9.17~\cite{CellListMap} for the
neighbour list. The tests in \texttt{test/runtests.jl} compare the field against the direct
periodic sum, and the derivative of $G_{\mathrm{near}}$ against finite differences. The numbers in
this report are produced by \texttt{papers/ewald-report/parameter\_study.jl}.

\begin{thebibliography}{99}

\bibitem{Ewald1921}
P.~P. Ewald, Die Berechnung optischer und elektrostatischer Gitterpotentiale,
\emph{Ann.\ Phys.} \textbf{369}, 253--287 (1921).

\bibitem{SalinCaillol2000}
G.~Salin and J.-M. Caillol, Ewald sums for Yukawa potentials,
\emph{J.\ Chem.\ Phys.} \textbf{113}, 10459--10463 (2000).

\bibitem{Darden1993}
T.~Darden, D.~York and L.~Pedersen, Particle mesh Ewald: An $N\log(N)$ method for Ewald sums in
large systems, \emph{J.\ Chem.\ Phys.} \textbf{98}, 10089--10092 (1993).

\bibitem{Essmann1995}
U.~Essmann, L.~Perera, M.~L. Berkowitz, T.~Darden, H.~Lee and L.~G. Pedersen, A smooth particle
mesh Ewald method, \emph{J.\ Chem.\ Phys.} \textbf{103}, 8577--8593 (1995).

\bibitem{LindboTornberg2011}
D.~Lindbo and A.-K. Tornberg, Spectral accuracy in fast Ewald-based methods for particle
simulations, \emph{J.\ Comput.\ Phys.} \textbf{230}, 8744--8761 (2011).

\bibitem{Barnett2019}
A.~H. Barnett, J.~Magland and L.~af~Klinteberg, A parallel nonuniform fast Fourier transform
library based on an ``exponential of semicircle'' kernel, \emph{SIAM J.\ Sci.\ Comput.}
\textbf{41}, C479--C504 (2019).

\bibitem{NonuniformFFTs}
J.~I. Polanco, NonuniformFFTs.jl, \url{https://github.com/jipolanco/NonuniformFFTs.jl}.

\bibitem{FFTW}
M.~Frigo and S.~G. Johnson, The design and implementation of FFTW3, \emph{Proc.\ IEEE}
\textbf{93}, 216--231 (2005).

\bibitem{Unser1993}
M.~Unser, A.~Aldroubi and M.~Eden, B-spline signal processing: Part I---Theory,
\emph{IEEE Trans.\ Signal Process.} \textbf{41}, 821--833 (1993).

\bibitem{Interpolations}
Interpolations.jl, \url{https://github.com/JuliaMath/Interpolations.jl}.

\bibitem{CellListMap}
L.~Martínez, CellListMap.jl: Efficient and customizable cell list implementation for calculation
of pairwise particle properties within a cutoff, \emph{Comput.\ Phys.\ Commun.} \textbf{279},
108452 (2022).

\end{thebibliography}

\end{document}
